% Capítulo 3, versão 3.
% Cada tecnologia segue a mesma ordem: o que é (com fonte), o que oferece e como é usada neste trabalho.

\chapter{Suporte Tecnológico}
\label{cap-suporte-tecnologico}

Este capítulo apresenta as ferramentas utilizadas ou previstas para a redação da monografia, a modelagem da arquitetura e o desenvolvimento da prova de conceito. Cada ferramenta é descrita na mesma ordem. Primeiro se apresenta o que ela é, depois os recursos que interessam ao trabalho e, por fim, o uso previsto. A Seção~\ref{st-tcc} trata das ferramentas de apoio geral, a Seção~\ref{st-modelagem}, das ferramentas de modelagem, e a Seção~\ref{st-desenvolvimento}, das tecnologias de implementação, comunicação, análise de software e interpretação documental. A Seção~\ref{st-resumo} resume o capítulo em um quadro com a finalidade e a situação de adoção de cada componente.

\section{Ferramentas de Apoio Geral}
\label{st-tcc}

\subsection{LaTeX}
\label{st-3-1-1}

LaTeX é um sistema de preparação de documentos que oferece recursos para composição de textos científicos, referências cruzadas, figuras, tabelas e bibliografias \cite{latexProject}. Neste trabalho, é utilizado para organizar e formatar a monografia a partir do modelo adotado pelo curso. Os arquivos dos capítulos, as referências bibliográficas e as figuras são reunidos nesse projeto de redação.

\subsection{Zotero}
\label{st-3-1-2}

Zotero é um gerenciador de referências que permite reunir registros bibliográficos, organizar fontes em coleções e associar arquivos e notas aos itens cadastrados \cite{zoteroGuide}. Neste trabalho, apoia a organização do acervo consultado e a recuperação dos documentos usados na fundamentação.

Os registros bibliográficos são conferidos quanto à autoria, ao título, ao ano e à identificação da publicação. A planilha de levantamento permanece como registro dos critérios e das decisões de seleção e distingue os estudos selecionados da literatura complementar.

\section{Ferramentas de Modelagem}
\label{st-modelagem}

\subsection{UML, draw.io e PlantUML}
\label{st-3-2-1}

A arquitetura é modelada em UML, linguagem gráfica especificada pelo \textit{Object Management Group} para visualizar, especificar, construir e documentar sistemas \cite{omgUml}. São usados dois diagramas. O diagrama de componentes mostra os agentes, as ferramentas e os artefatos, com suas relações. O diagrama de sequência mostra a ordem das mensagens em uma auditoria.

O diagrama de componentes é elaborado no draw.io, também conhecido como diagrams.net. A ferramenta conserva os diagramas em formato editável e exporta figuras em formatos como SVG e PDF \cite{drawioExport}. O diagrama de sequência é mantido em PlantUML, ferramenta em que o diagrama é definido inteiramente em texto e desenhado automaticamente \cite{plantumlSequence}.

Os diagramas distinguem a arquitetura alvo do recorte implementado no protótipo inicial. Os arquivos editáveis são conservados para acompanhar a evolução do projeto, e as figuras são apresentadas e discutidas no capítulo da Proposta.

\section{Ferramentas de Desenvolvimento}
\label{st-desenvolvimento}

\subsection{Python}
\label{st-3-3-1}

Python é a linguagem de implementação dos agentes, dos adaptadores das ferramentas e do processamento dos registros da auditoria. Sua biblioteca \texttt{asyncio} oferece recursos de programação assíncrona, comunicação de rede e controle de subprocessos \cite{pythonAsyncio}, necessários ao envio de mensagens e ao acompanhamento das ferramentas externas.

A escolha da linguagem decorre do uso de SPADE, apresentado a seguir. O protótipo inicial registra o uso de Python 3.13.7. As versões efetivamente utilizadas na aplicação da pesquisa serão conservadas junto aos registros de execução.

A linguagem de implementação é independente das linguagens dos repositórios analisados. A cobertura de Python, TypeScript e PHP depende das ferramentas, dos artefatos disponíveis e da configuração de cada análise. A demonstração inicial utiliza um repositório sintético em Python, anterior à aplicação em repositórios autorizados da Empresa A.

\subsection{SPADE}
\label{st-3-3-2}

SPADE, sigla de \textit{Smart Python Agent Development Environment}, é uma plataforma de sistemas multiagentes escrita em Python e baseada em XMPP \cite{palanca2020,spadeDocs}. Sua documentação descreve um modelo de agentes organizado por comportamentos e apoiado em \texttt{asyncio}, e informa que a plataforma pode ser usada com qualquer servidor XMPP \cite{spadeDocs}.

Na prova de conceito, SPADE é empregado na construção dos agentes e no envio e recebimento das mensagens. O protótipo inicial utiliza a versão 4.1.4 e mantém quatro participantes em processos separados, que são o Orquestrador, o agente de Qualidade, o de Segurança e o Auditor. O agente de Requisitos e Documentação integra a arquitetura alvo e será acrescentado na etapa de interpretação documental.

Essa separação conserva a identidade dos participantes e permite acompanhar a origem das mensagens. As consultas P2P são encaminhadas entre os especialistas. O Orquestrador organiza a execução, e o Auditor recebe as conclusões para consolidação. SPADE fornece a infraestrutura de agentes, e os comportamentos e os critérios de análise são especificados pelo projeto.

\subsection{XMPP e Prosody}
\label{st-3-3-3}

XMPP é um protocolo que permite a troca, quase em tempo real, de dados estruturados entre duas ou mais entidades de rede \cite{rfc6120}. É o protocolo de transporte utilizado por SPADE. Prosody é o servidor XMPP selecionado para o ambiente inicial, com recursos documentados de configuração, módulos e certificados \cite{prosodyDocs}. Como o protótipo usa uma compilação de desenvolvimento do servidor, o pacote ou a imagem utilizados serão identificados nos registros do ambiente, para permitir a reprodução da configuração. Cada agente utiliza uma identidade própria na comunicação.

A comunicação P2P adotada é a relação lógica entre os agentes que formulam e respondem às consultas. O servidor apenas transporta as mensagens. Uma solicitação do agente de Segurança continua sendo respondida diretamente pelo agente de Qualidade ou pelo de Requisitos e Documentação.

O protótipo registra os identificadores das auditorias e das solicitações, para relacionar perguntas, respostas e conclusões. Tempos de espera e estados de falha permitem identificar interações não concluídas. Esses registros serão utilizados na análise do funcionamento da arquitetura, conforme os procedimentos definidos na Metodologia.

\subsection{Docker Compose}
\label{st-3-3-4}

Docker Compose permite definir e executar aplicações compostas por múltiplos contêineres, reunindo configurações de serviços, redes e volumes em um arquivo YAML \cite{dockerCompose}. No protótipo, esse recurso organiza os serviços de infraestrutura necessários à comunicação e à análise de código.

As configurações do ambiente são conservadas para identificar os serviços utilizados e suas versões. As credenciais de acesso são fornecidas por configuração externa ao código-fonte, incluindo variáveis de ambiente para os \textit{tokens} exigidos pelas ferramentas. Assim, os valores de acesso não integram os relatórios da pesquisa.

\subsection{SonarQube e SonarScanner}
\label{st-3-3-5}

SonarQube produz medidas de qualidade e aponta violações das regras de análise. O SonarScanner executa a análise configurada e envia um relatório ao servidor, que o processa em uma tarefa de segundo plano. Segundo a documentação, uma análise só está completa quando essa tarefa termina \cite{sonarAnalysis99}.

No protótipo inicial, SonarQube fornece as evidências consultadas pelo agente de Qualidade. O adaptador acompanha o processamento da tarefa e relaciona os resultados à análise e à revisão do repositório. Assim, o fim da execução do \textit{scanner} não é tomado como sinal de que os resultados já estão disponíveis.

As medidas conservam sua granularidade. Um resultado agregado por arquivo é interpretado nesse nível, sem atribuição automática a uma função específica. O estado do \textit{Quality Gate} é conservado como informação da ferramenta. O parecer o apresenta sem convertê-lo em veredito, e a classificação cabe ao responsável humano.

O ambiente inicial registra SonarQube Community 9.9.8.100196 e SonarScanner CLI 8.1.0.6389. A cobertura efetiva dos repositórios selecionados será verificada antes da aplicação, pois depende da linguagem e da edição da ferramenta.

\subsection{Semgrep}
\label{st-3-3-6}

Semgrep é uma ferramenta de análise estática que identifica padrões de código por meio de regras e produz um relatório dos achados. Sua edição comunitária analisa Python, TypeScript, JavaScript e PHP. A documentação classifica essas linguagens, nessa edição, como suportadas pela comunidade, sem prioridade nas atualizações de regras e do motor, e limita a análise ao interior de cada função \cite{semgrepDocs}. Essas duas características serão registradas como limitações da cobertura.

O Semgrep será a fonte de evidências de segurança estática para o agente de Segurança. Ele substitui o Bandit, ferramenta restrita a Python utilizada no protótipo inicial. O uso de uma única ferramenta para as três linguagens evita adaptadores distintos e uniformiza o formato das evidências. Os resultados preservam a identificação da regra, a localização e as classificações emitidas pela ferramenta. O especialista pode consultar medidas do SonarQube referentes ao mesmo arquivo para contextualizar um achado, conservada a distinção entre complexidade e evidência de vulnerabilidade.

O motor da edição comunitária é distribuído sob a licença LGPL 2.1, e as regras mantidas pelo fabricante são cobertas por uma licença que restringe seu uso a fins internos \cite{semgrepLicensing}. Esse uso será considerado nas condições de aplicação da pesquisa. As regras serão mantidas em cópia local versionada, e o conjunto de regras e sua versão serão registrados em cada análise, pois a cobertura depende das regras disponíveis para cada linguagem. O protótipo não executa a aplicação auditada nem modifica seus arquivos.

\subsection{Ferramentas para Dependências e Ampliação de Cobertura}
\label{st-3-3-7}

\begin{sloppypar}
A análise de dependências identifica os componentes e as versões que o software utiliza e os relaciona a vulnerabilidades conhecidas \cite{blackIr8397}. Essa integração ainda não compõe o protótipo inicial. As ferramentas do Quadro~\ref{quadro-ferramentas-dependencias} são candidatas e serão selecionadas conforme os manifestos e os gerenciadores encontrados nos repositórios.
\end{sloppypar}

\begin{table}[p]
\centering
\caption{Ferramentas candidatas à análise de dependências e ampliação de cobertura}
\label{quadro-ferramentas-dependencias}
\begingroup
\footnotesize
\SingleSpacing
\setlength{\tabcolsep}{4pt}
\renewcommand{\arraystretch}{1.15}
\begin{tabularx}{\textwidth}{>{\raggedright\arraybackslash}X>{\raggedright\arraybackslash}X>{\raggedright\arraybackslash}X}
\hline
\textbf{Ferramenta} & \textbf{Aplicação prevista} & \textbf{Condição de utilização} \\ \hline
pip-audit & Consulta de vulnerabilidades conhecidas em dependências Python & Identificação das dependências e versões em formato compatível \\ \hline
npm audit & Auditoria de dependências de projetos que utilizam npm, inclusive aplicações TypeScript & Presença de estrutura de dependências compatível com npm \\ \hline
Composer audit & Consulta de avisos de segurança associados a pacotes gerenciados por Composer & Projeto PHP que utilize Composer e permita identificar as versões analisadas \\ \hline
Trivy & Verificações adicionais de vulnerabilidades, segredos e configurações, conforme alvo e modo & Seleção explícita do escopo e dos formatos suportados \\ \hline
\end{tabularx}
\endgroup
\legend{Fonte: elaboração do autor.}
\end{table}

\begin{sloppypar}
As funções dessas ferramentas são descritas em suas documentações oficiais \cite{pipAuditDocs,npmAuditDocs,composerAuditDocs,trivyDocs}. Relatórios sobrepostos de ferramentas diferentes não serão tratados como evidências independentes de um mesmo problema. A data da consulta e a identificação das bases de vulnerabilidades utilizadas serão conservadas.
\end{sloppypar}

Os adaptadores são configurados apenas para coleta de evidências. Comandos de correção automática, de atualização de pacotes ou de execução da aplicação não integram o procedimento. A ausência de um manifesto compatível será registrada como limitação de cobertura, e não como ausência de vulnerabilidades.

\subsection{Modelo de Linguagem e Execução Local}
\label{st-3-3-8}

O agente de Requisitos e Documentação utiliza um modelo de linguagem para interpretar documentos textuais a partir das consultas dos especialistas. Como em \citeonline{yesil2026}, a extração de informações dos documentos fica separada da decisão orientada por regras. O modelo fornece contexto. A sinalização de pontos de atenção permanece com as regras do Auditor, e a decisão final, com o responsável humano.

A execução local do modelo está em avaliação. Entre os recursos considerados está o Ollama, que oferece uma interface de programação para integrar aplicações a modelos executados localmente \cite{ollamaApi}. Nessa configuração, o agente encaminha as consultas ao serviço de inferência e recebe as respostas para processamento e registro. A execução local mantém o processamento dos documentos no ambiente definido para a pesquisa, observadas as condições de acesso e de tratamento dos dados da organização.

A seleção do modelo considerará a capacidade de interpretar documentos em português, a fundamentação das respostas nos trechos fornecidos, o reconhecimento de informação insuficiente, o tempo de resposta e o consumo de memória. O tamanho do contexto também será considerado, pois um contexto maior aumenta a memória necessária para executar o modelo \cite{ollamaContext}. A adoção do Ollama e a escolha do modelo dependem da verificação de sua adequação à tarefa.

Serão registrados a identificação e a versão do modelo, a quantização, quando utilizada, o tamanho de contexto e os parâmetros de geração, juntamente com as entradas e as respostas produzidas. A integração deverá receber a consulta do especialista, identificar os documentos pertinentes e conservar os trechos utilizados na resposta. Deverá também registrar informação insuficiente, resposta sem apoio documental e conflito entre fontes, estados que distinguem um contexto recuperado de uma conclusão que não pôde ser sustentada.

A recuperação aumentada por geração, empregada por \citeonline{yesil2026}, é uma possibilidade de implementação. Sua adoção, o mecanismo de busca e a eventual indexação serão especificados após a definição do conjunto de documentos, e não há banco vetorial escolhido nesta etapa. Os critérios de aceitação das respostas e a ligação entre arquivos e componentes são detalhados no capítulo da Proposta.

\subsection{Versionamento, Registros e Apresentação}
\label{st-3-3-9}

Git é empregado na identificação das revisões dos repositórios examinados e no acompanhamento das alterações do projeto. O sistema conserva o histórico dos repositórios e permite relacionar o código a uma revisão identificada \cite{gitAbout}. Na auditoria, o identificador do \textit{commit} delimita a versão analisada.

O protótipo utiliza arquivos JSON para evidências e pareceres e registros JSONL para eventos da execução. As regras do Auditor são armazenadas em YAML, o que permite conservar a política utilizada em cada auditoria. As dependências do código incluem Pydantic, HTTPX e PyYAML, usadas, respectivamente, na validação estrutural dos dados, nas consultas HTTP às ferramentas e no carregamento das políticas.

A identificação do \textit{commit}, das ferramentas e das configurações relaciona o parecer à versão examinada. Os resultados originais das ferramentas são conservados junto às evidências normalizadas e às mensagens pertinentes. Assim, o relatório permite localizar a fonte de uma afirmação e acompanhar sua participação na conclusão.

Para a apreciação humana estruturada, será preparado um material de leitura com o parecer, as evidências citadas, os trechos documentais e as regras acionadas. O responsável humano avaliará a clareza, a coerência e a fundamentação desse material, sem realizar uma nova auditoria. O instrumento de registro das apreciações é definido na Metodologia.

A apresentação consolidada das auditorias poderá incluir um painel, acompanhado dos relatórios detalhados. A tecnologia dessa interface ainda não foi selecionada. Ela deverá preservar o acesso às evidências e às limitações de cada execução, para que a visão agregada não oculte diferenças de cobertura entre repositórios.

\section{Resumo do Capítulo}
\label{st-resumo}

O suporte tecnológico reúne recursos de redação, de modelagem e de implementação. Python e SPADE sustentam os agentes, e Prosody transporta as mensagens por XMPP. SonarQube fornece as evidências de qualidade, e o Semgrep, as de segurança estática. A interpretação documental e a análise de dependências serão integradas nas etapas seguintes, conforme a seleção dos componentes e o escopo dos repositórios. O Quadro~\ref{quadro-tecnologias} resume a finalidade e a situação de cada tecnologia.

\begin{table}[p]
\centering
\caption{Tecnologias e recursos utilizados ou previstos no trabalho}
\label{quadro-tecnologias}
\begingroup
\footnotesize
\SingleSpacing
\setlength{\tabcolsep}{4pt}
\renewcommand{\arraystretch}{1.15}
\begin{tabularx}{\textwidth}{>{\raggedright\arraybackslash}X>{\raggedright\arraybackslash}X>{\raggedright\arraybackslash}X>{\raggedright\arraybackslash}X}
\hline
\textbf{Tecnologia ou recurso} & \textbf{Categoria} & \textbf{Finalidade no trabalho} & \textbf{Versão ou situação} \\ \hline
LaTeX & Apoio geral & Redação e formatação da monografia & Projeto de redação existente; distribuição a registrar \\ \hline
Zotero & Apoio geral & Organização do acervo e das referências & Em uso; versão a registrar \\ \hline
UML, draw.io e PlantUML & Modelagem & Diagramas de componentes e de sequência & Em uso; versões a registrar \\ \hline
Python & Desenvolvimento & Implementação dos agentes e adaptadores & 3.13.7 no protótipo inicial \\ \hline
SPADE & Desenvolvimento & Comportamentos e mensagens dos agentes & 4.1.4 no protótipo inicial \\ \hline
Prosody & Desenvolvimento & Servidor XMPP & 13.0 nightly build 100 no protótipo inicial \\ \hline
Docker Compose & Desenvolvimento & Configuração dos serviços & Utilizado; versão a registrar \\ \hline
SonarQube Community & Desenvolvimento & Medidas e achados de código & 9.9.8.100196 no protótipo inicial \\ \hline
SonarScanner CLI & Desenvolvimento & Envio das análises ao SonarQube & 8.1.0.6389 no protótipo inicial \\ \hline
Semgrep (edição comunitária) & Desenvolvimento & Análise estática de segurança em Python, TypeScript e PHP & Integração pendente; versão a registrar. O protótipo inicial utilizou Bandit 1.9.4 \\ \hline
pip-audit, npm audit, Composer audit e Trivy & Desenvolvimento & Análise de dependências e ampliação da cobertura & Candidatos; seleção e integração pendentes \\ \hline
Ollama e modelo de linguagem & Desenvolvimento & Serviço de interpretação documental & Em avaliação; modelo a selecionar \\ \hline
Git & Desenvolvimento & Histórico e identificação das revisões & Utilizado; versão a registrar \\ \hline
JSON, JSONL e YAML & Formatos de apoio à implementação & Evidências, eventos e políticas & Utilizados no protótipo inicial \\ \hline
Painel e material de leitura dos pareceres & Apresentação & Consulta aos resultados e às evidências & Tecnologia da interface e instrumento a definir \\ \hline
\end{tabularx}
\endgroup
\legend{Fonte: elaboração do autor, a partir do projeto e dos registros do protótipo inicial.}
\end{table}

As versões indicadas documentam o ambiente já empregado e serão atualizadas conforme a configuração efetiva da aplicação da pesquisa. Os procedimentos de avaliação e apreciação humana são definidos na Metodologia. A arquitetura, os diagramas e os critérios de consolidação dos pareceres são detalhados no capítulo da Proposta.
